\chapter{Anexo 1 - Implicaciones sociales y éticas del proyecto}

%En un proyecto de esta índole, se ven muchos factores incluidos. La vida humana, la calidad de la vida humana, el nombre de una empresa, el nombre de un ingeniero, la conciencia de un ser humano, el precio de la vida y la importancia de la responsabilidad que lleva realizar un trabajo que pueda afectar una.
%En un proceso el cual tiene cómo objetivo crear un producto del cual dependen vidas no es algo simple, o aleatorio. Es un proceso importante, delicado y digno de recibir toda la atención que amerita la dependencia de la vida de un ser humano en el resultado de este proceso.
%¿Que es más importante? ¿El dinero? ¿La productividad? o ¿la vida humana?
%En repetidas ocasiones se han visto decisiones tomadas por corporaciones y entes que terminan por indicar que para estos los primeros dos eran la respuesta.
%Con mi proyecto se buscaba demostrar la importancia de la ultima, de la vida humana. Un proyecto con el objetivo de mejorar la forma en la que se determina la calidad de un producto y al mismo tiempo mejorar la ganancia monetaria. Porque al fin y al cabo ambos son un objetivo, sólo que el segundo nunca debería opacar al primero.
%Un error de mi parte, podría significar la muerte de una persona, lo cual no es una responsabilidad menor, y con toda la responsabilidad profesional y todo el deber humano de cuidar los unos de los otros debe ser tratada.
%Mi deber como ingeniero, es cumplir con mi trabajo con la mayor de mis capacidades y con la mayor responsabilidad. Mi conocimiento está al servicio de la humanidad.

%Un factor, también importante, que muchas veces es omitido es el factor de la automatización sobre el proceso humano. Automatizar un proceso, volver más efectivo y confiable que con las personas aunque sea una necesidad y un objetivo imperativo, tiene consecuencias posibles. Ya que al retirar a un humano de su puesto para ser reemplazado por una maquina se causa que alguien pierda su ingreso monetario, y pueda terminar en una situación de desdicha y necesidad. Para mi tranquilidad, mi proyecto aunque busque reducir el componente humano no lo elimina completamente, pues aquellos que lo utilizarán no verán su trabajo retirado si no que simplificado.

%Para concluir esta reflexión, deseo hacer énfasis en la necesidad del progreso, en la necesidad de la mejora a base de la automatización, pero al mismo tiempo hacer énfasis en la humanidad de nuestra profesión, en la responsabilidad moral y profesional que tenemos con nuestros semejantes, y en la belleza de aquello que hacemos para ayudarlos.

